\section{Transfert de Markov compatible avec toute la limite projective}
\label{sec:cpe-transferencia-limite-inverso}

Soit \(\rho_m:H_m\twoheadrightarrow Y_m\) le quotient observable à la profondeur
\(m\) et soit \(L_m(y,y')\) le noyau de transition déjà construit dans
\(Y_m\). Pour chaque fibre finie \(\rho_m^{-1}(y)\), soit
\(\nu_{m,y}\) la mesure uniforme conditionnelle. La régularité nonadique des
extensions permet de choisir ces mesures de sorte que leur image par la
troncature soit \(\nu_{m-1,\bar y}\).

On définit, pour \(h\in H_m\),
\begin{equation}
 \mathcal K_m(h,A)
 :=\sum_{y'\in Y_m}
 L_m(\rho_m(h),y')\,
 \nu_{m,y'}\bigl(A\cap\rho_m^{-1}(y')\bigr).
 \label{eq:cpe-kernel-lift-uniforme}
\end{equation}

\begin{theorem}[Relèvement fortement agrégeable et limite de Markov]
\label{thm:cpe-kernel-proyectivo}
Les noyaux \(\mathcal K_m\) satisfont
\begin{equation}
 \sum_{\rho_m(g)=y'}\mathcal K_m(h,\{g\})
 =L_m(\rho_m(h),y'),
 \label{eq:cpe-agregabilidad-fuerte}
\end{equation}
indépendamment du représentant \(h\) de la fibre. De plus, ils commutent avec
les troncatures. Il existe donc un unique noyau de Markov
\(\mathcal K_\infty\) sur la limite projective dont la marginale de niveau \(m\)
est \(\mathcal K_m\). Sa projection par \(\rho_\infty\) est le noyau
observable \(L_\infty\).
\end{theorem}

\begin{proof}
La somme du membre de gauche de
\eqref{eq:cpe-agregabilidad-fuerte} est
\[
 L_m(\rho_m(h),y')
 \sum_{g\in\rho_m^{-1}(y')}\nu_{m,y'}(\{g\})
 =L_m(\rho_m(h),y').
\]
La compatibilité de \(L_m\), de \(\rho_m\) et des mesures conditionnelles
prouve le carré de troncature. Les distributions cylindriques obtenues
en itérant \(\mathcal K_m\) sont cohérentes ; le théorème d'extension pour les
systèmes projectifs d'espaces finis produit \(\mathcal K_\infty\), et
l'unicité découle du fait que les cylindres engendrent la sigma-algèbre de la limite.
\end{proof}

La construction conserve l'immobilité microscopique et les étiquettes complètes
dans la mesure conditionnelle. La moyenne visible ne remplace pas l'état
enrichi et ne s'identifie pas au Topological Phase Kernel.

